\documentclass[11pt,aspectratio=169]{beamer}
\usetheme{Boadilla}
\usefonttheme{professionalfonts}
\setbeamertemplate{navigation symbols}{}

\usepackage[T1]{fontenc}
\usepackage{lmodern}
\usepackage{amsmath,amssymb,mathtools}
\usepackage{booktabs}
\usepackage{tabularx}
\usepackage{array}
\usepackage[onehalfspacing]{setspace}

\newcommand{\figplaceholder}[1]{\begin{center}\fbox{\begin{minipage}[c][0.32\textheight][c]{0.82\textwidth}\centering\large #1\end{minipage}}\end{center}}
\newcommand{\shortfigplaceholder}[1]{\begin{center}\fbox{\begin{minipage}[c][0.22\textheight][c]{0.82\textwidth}\centering\large #1\end{minipage}}\end{center}}
\newcommand{\compactfigplaceholder}[1]{\begin{center}\fbox{\begin{minipage}[c][0.10\textheight][c]{0.82\textwidth}\centering\normalsize #1\end{minipage}}\end{center}}
\newcommand{\R}{\mathcal{R}}
\newcommand{\E}{\mathcal{E}}
\newcommand{\dd}{\,\mathrm{d}}

\title[Tidal Lanes]{Tidal Traffic}
\author[Duan (SJTU), Feng (PKU), Gu (PKU)]{Yige Duan (SJTU) $\quad$ Xiaoruo Feng (PKU) $\quad$ Yizhen Gu (PKU)}
\date{\today}

\begin{document}

% 1
\begin{frame}
\titlepage
\end{frame}

% 10
\begin{frame}{Motivation}
\begin{itemize}
  \item \textbf{Peak-direction imbalance}
  \begin{itemize}
    \item Urban roads often face strong directional congestion
    \item Commuting traffic is shaped by where workers live and where jobs are located
  \end{itemize}
  \item Daily two-way cost
  \begin{itemize}
    \item Return trips reverse the same flows
    \item Value-of-time costs should account for both commute and return trips
  \end{itemize}
  \item Policy margin
  \begin{itemize}
    \item Tidal lanes reallocate existing lanes across directions
    \item They may change route choices and residence--workplace sorting
  \end{itemize}
\end{itemize}
\end{frame}

% 11
\begin{frame}{What We Do}
\begin{itemize}
  \item Model \textcolor{blue}{directional congestion}
  \begin{itemize}
    \item Extend the Allen--Arkolakis inner-city model
    \item Keep physical road infrastructure symmetric for each road segment
  \end{itemize}
  \item Generate asymmetric commuting costs
  \begin{itemize}
    \item Traffic per lane differs by direction
    \item Residence and workplace choices determine directional flows
  \end{itemize}
  \item Evaluate tidal lanes
  \begin{itemize}
    \item Hold total lanes fixed and shift one lane toward the more congested direction
    \item Decompose local congestion, global congestion, and residence--workplace sorting
  \end{itemize}
\end{itemize}
\end{frame}

% 12
\begin{frame}{Literature}
\begin{itemize}
  \item Spatial equilibrium and agglomeration
  \begin{itemize}
    \item Allen and Arkolakis (RES, 2022)
    \item Ahlfeldt et al. (2015), as used by Allen--Arkolakis
  \end{itemize}
  \item Transportation policy and congestion
  \begin{itemize}
    \item Barwick et al. (AER, 2024)
    \item Yang et al. (AEJ: Policy, 2020)
  \end{itemize}
  \item Calibration and contribution
  \begin{itemize}
    \item Beijing studies estimate \textcolor{blue}{$\delta_1=-1.1$}
    \item This project adds a \underline{tidal-lane margin}
  \end{itemize}
\end{itemize}
\end{frame}

% 13
\begin{frame}{Data}
\small
\begin{itemize}
  \item Network and traffic
  \begin{itemize}
    \item Road endpoints, length, direction, lanes, and functional class
    \item Link-level traffic volume, ideally by direction and time of day
  \end{itemize}
  \item Speeds
  \begin{itemize}
    \item Observed speed or travel time by segment and time interval
    \item Summary objects for Figure 8 and Table 9
  \end{itemize}
  \item Commuting and geography
  \begin{itemize}
    \item Residence location, workplace location, and OD commuting flows
    \item Grid cells, tract boundaries, and CBD or subcenter indicators
  \end{itemize}
\end{itemize}
\compactfigplaceholder{Figure 1 placeholder: study area and main road network map}
\end{frame}

% 14
\begin{frame}{Road Segments}
\small
\begin{itemize}
  \item Build model links
  \begin{itemize}
    \item Convert raw road geometry into links between adjacent nodes
    \item Clean duplicate or unusable road records
  \end{itemize}
  \item Preserve direction-specific attributes
  \begin{itemize}
    \item Keep lanes and peak traffic where available
    \item Treat each physical road as symmetric infrastructure with two directed flows
  \end{itemize}
  \item Document network construction
  \begin{itemize}
    \item Figure 2: road segments
    \item Figure 4: selected or cleaned network
  \end{itemize}
\end{itemize}
\begin{columns}[T,onlytextwidth]
\column{0.49\textwidth}
\compactfigplaceholder{Figure 2 placeholder: road segments}
\column{0.49\textwidth}
\compactfigplaceholder{Figure 4 placeholder: selected or cleaned network}
\end{columns}
\end{frame}

% 15
\begin{frame}{Speed Statistics}
\small
\begin{equation}
v_{klh}=\frac{\mathrm{distance}_{kl}}{\mathrm{travel\ time}_{klh}},
\qquad
m_{klh}=v_{klh}^{-1}.
\end{equation}
\begin{itemize}
  \item Speed and inverse speed
  \begin{itemize}
    \item $v_{klh}$ is speed on road segment $k$--$l$ in time interval $h$
    \item $m_{klh}$ is \textbf{inverse speed}, interpreted as minutes per unit distance
  \end{itemize}
  \item Congestion summaries
  \begin{itemize}
    \item Figure 8 shows the speed distribution
    \item Table 9 reports speed summary statistics
  \end{itemize}
\end{itemize}
\begin{columns}[T,onlytextwidth]
\column{0.55\textwidth}
\compactfigplaceholder{Figure 8 placeholder: speed distribution}
\column{0.43\textwidth}
\compactfigplaceholder{Table 9 placeholder: speed summary statistics}
\end{columns}
\end{frame}

% 16
\begin{frame}{Speed Within A Day}
\small
\begin{itemize}
  \item Peak timing
  \begin{itemize}
    \item Compare speeds across time intervals
    \item Identify morning and evening peaks
  \end{itemize}
  \item Directional imbalance
  \begin{itemize}
    \item Compare opposite directions within the same road segment
    \item Use within-day patterns to reveal \textcolor{red}{tidal congestion}
  \end{itemize}
  \item Model discipline
  \begin{itemize}
    \item Map traffic per lane into speed
    \item Calibrate the speed elasticity from the literature
  \end{itemize}
\end{itemize}
\compactfigplaceholder{Figure 9 placeholder: within-day speed profile}
\end{frame}

% 17
\begin{frame}{Tidal Commuting}
\small
\begin{itemize}
  \item \textbf{Tidal commuting}
  \begin{itemize}
    \item Dominant commute direction reverses over the day
    \item Morning and return-trip asymmetry are two sides of the same flow pattern
  \end{itemize}
  \item Residence--workplace separation
  \begin{itemize}
    \item More workers may live in one area and work in another
    \item The imbalance generates asymmetric traffic per lane
  \end{itemize}
  \item Empirical object
  \begin{itemize}
    \item Measure \textcolor{blue}{directional traffic per lane}
    \item Use Figure 10 to summarize tidal commuting patterns
  \end{itemize}
\end{itemize}
\compactfigplaceholder{Figure 10 placeholder: tidal commuting patterns}
\end{frame}

% 18
\begin{frame}{Grids}
\small
\begin{itemize}
  \item Spatial aggregation
  \begin{itemize}
    \item Aggregate locations into grid cells
    \item Use cells for residence, workplace, and network nodes
  \end{itemize}
  \item Data linkage
  \begin{itemize}
    \item Link OD commuting data to the road network
    \item Preserve consistent geography across model objects
  \end{itemize}
  \item Counterfactual units
  \begin{itemize}
    \item Reallocate residence and workplace shares across grid cells
    \item Use Figure 14 to document the analysis grid
  \end{itemize}
\end{itemize}
\compactfigplaceholder{Figure 14 placeholder: analysis grid}
\end{frame}

% 2
\begin{frame}{Model: Environment}
\begin{itemize}
  \item Locations $i,j\in N$ host $\bar{L}$ mass of residents and jobs
  \vspace{0.5em}
  \item A physical road is represented by two directed links, $(k,l)$ and $(l,k)$
  \begin{itemize}
    \item The road has symmetric quality $t_{kl}=t_{lk}$, but directed lanes $n_{kl}$ and $n_{lk}$ can differ
    \item A tidal-lane policy shifts one lane across directions while holding $n_{kl}+n_{lk}$ fixed
  \end{itemize}
  \vspace{0.5em}
  \item Traversal cost is \textcolor{blue}{endogenous} and \textcolor{red}{directed}
  \begin{itemize}
    \item $c_{kl}$ depends on road quality and traffic-per lane
    \item $c_{kl} \neq c_{lk}$ when traffic differs by direction
  \end{itemize}
\end{itemize}
\end{frame}

% 3
\begin{frame}{Individual Location Choice}
\begin{itemize}
  \item An individual $\omega$ living in $i$, working in $j$, and taking route $r\in\R_{ij}$ obtains
  \vspace{-0.5em}\begin{footnotesize}\begin{equation}
    V_{ij,r}(\omega) = \left( \frac{u_i w_j}{C_{ij,r}} \right)\varepsilon_{ij,r}(\omega),
    \qquad
    C_{ij,r} = \prod_{\ell=1}^{K_r} c_{r_{\ell-1},r_\ell}
  \vspace{-1.25em}\end{equation}\end{footnotesize}
  \begin{itemize}
    \item $u_i$ is amenity; $w_j$ is wage; $\epsilon$ is i.i.d. Frechet
    \item Route $r = [r_0, r_1, \cdots, r_{K_r}]$ with $r_0=i$ and $r_{K_r}=j$
  \end{itemize}
  \vspace{0.5em}
  \item Choice probabilities are
  \vspace{-0.5em}\begin{equation}\label{eq:pi}
    \pi_{ij,r} = \frac{\left(u_i w_j/C_{ij,r}\right)^\theta}{\sum_{a}\sum_b\sum_{s\in\R_{ab}}\left(u_a w_b/C_{ab,s}\right)^\theta}
  \end{equation}
\end{itemize}
\end{frame}

% 4
\begin{frame}{Bilateral Commuting Costs and Flows}
\begin{itemize}
  \item The bilateral commuting-cost, aggregated across all routes (Frechet expectation)
  \vspace{-0.5em}\begin{equation}
    \tau_{ij} = \left[ \sum_{r\in\R_{ij}} C_{ij,r}^{-\theta} \right]^{-1/\theta}
  \end{equation}
  \item Thus, residence-workplace choice shares satisfy
  \vspace{-0.25em}\begin{equation}\label{eq:gravity}
    L_{ij} = \sum_{r\in\R_{ij}}L_{ij,r} = \tau_{ij}^{-\theta} \left(\frac{L_i^R}{\Pi_i}\right) \left(\frac{L_j^F}{P_j}\right), \,\, L_i^R=\sum_jL_{ij}, \,\, L_j^F=\sum_iL_{ij}
  \vspace{-1em}\end{equation}
  \begin{itemize}
    \item Directional congestion implies $\tau_{ij}\neq\tau_{ji}$ in general
    \item $\Pi_i$ and $P_j$ are origin and destination market-access terms
  \end{itemize}
\end{itemize}
\end{frame}

% 5
\begin{frame}{Amenity, Productivity, and Location Choices}
\begin{itemize}
  \item Agglomeration of local fundamentals
  \vspace{-0.5em}\begin{equation}\label{eq:agglom}
    w_j = A_j = \bar A_j\left(L_j^F\right)^\alpha, \qquad u_i=\bar u_i\left(L_i^R\right)^\beta
  \end{equation}
  \item Resident ($R$) and workplace ($F$) distribution across space
  \begin{align}
  \left( \frac{L_i^R}{\bar{L}}\right)^{1+\beta} & = \lambda \sum_j\tau_{ij}^{-\theta}\bar u_i\bar A_j \left(\frac{L_j^F}{\bar{L}}\right)^\alpha \label{eq:clear1} \\ 
  \left(\frac{L_i^F}{\bar{L}}\right)^{1+\alpha} & = \lambda \sum_j\tau_{ji}^{-\theta}\bar u_j\bar A_i \left(\frac{L_j^R}{\bar{L}}\right)^\beta \label{eq:clear2}
  \end{align}
  where $\lambda$ is a constant related to the expected welfare of residents
\end{itemize}
\end{frame}

% 34
\begin{frame}{Route-Specific Traversal Cost vs. Residence-Workplace Commuting Cost}
\begin{itemize}
  \item Define the directed adjacency matrix $A=[a_{kl}]$ by
  \vspace{-0.5em}\begin{equation}
    a_{kl} = c_{kl}^{-\theta}
  \end{equation}
  \item Summing over all feasible route lengths gives
  \vspace{-0.5em}\begin{equation}
    B = \sum_{K=0}^{\infty}A^K = (I-A)^{-1}
  \end{equation}
  \item The bilateral commuting-cost aggregator can be written as
  \vspace{-0.5em}\begin{equation}\label{eq:inverse}
    \tau_{ij} = B_{ij}^{-1/\theta}
  \end{equation}
\end{itemize}
\end{frame}

% 37
\begin{frame}{Feedback between Location Choice, Route Traffic, and Traversal Cost}
\begin{itemize}
  \item For a directed trip $i\to j$, define $\kappa^{kl}_{ij}$ as expected use of link $(k,l)$
  \vspace{-1em}\begin{equation}
    \kappa^{kl}_{ij} = \left( \frac{\tau_{ij}} {\tau_{ik}c_{kl}\tau_{lj}} \right)^\theta 
  \vspace{-0.5em}\end{equation}
  \item The directional traffic is
  \begin{equation}\label{eq:Xi_c}
  \Xi_{kl} = \sum_i\sum_j \kappa^{kl}_{ij}L_{ij}
  \end{equation}
  \item Traffic, in turn, affects route-level traversal cost
  \vspace{-0.5em}\begin{equation}\label{eq:c_Xi}
    c_{kl} = t_{kl} \left(\frac{\Xi_{kl}}{n_{kl}}\right)^{-\delta_1}
  \vspace{-0.5em}\end{equation}
    This allows for asymmetric traversal costs $c_{kl}\neq c_{lk}$ even if $t_{kl}=t_{lk}$
  \end{itemize}
\end{frame}

% 9
\begin{frame}{Equilibrium}
\begin{itemize}
  \item Given fundamentals $\{\bar u_i,\bar A_i,t_{kl},n_{kl}\}$ with $t_{kl}=t_{lk}$ on each physical road and parameters $(\theta,\alpha,\beta,\delta_1)$, an equilibrium is
  \vspace{-0.5em}\begin{equation}
    \left\{L_i^R,L_i^F,L_{ij},\tau_{ij},c_{kl},\Xi_{kl},\Pi_i,P_j,\lambda\right\}
  \vspace{-0.5em}\end{equation}
  such that
  \begin{itemize}
    \item residents choose residence-workplace-route options according to \eqref{eq:pi}
    \item bilateral flows satisfy commuting gravity \eqref{eq:gravity} and market clearing \eqref{eq:clear1}--\eqref{eq:clear2} 
    \item amenities and productivity satisfy \eqref{eq:agglom}
    \item bilateral costs are generated by the directed road network \eqref{eq:inverse}
    \item directed link traffic and directed traversal costs are mutually consistent \eqref{eq:Xi_c}--\eqref{eq:c_Xi}
  \end{itemize}
  \item Existence and uniqueness: see Allen and Arkolakis (2022)
\end{itemize}
\end{frame}

\begin{frame}{Exact-Hat Algebra for Tidal Lanes}
\begin{itemize}
  \item Exact-hat notation writes counterfactual changes as ratios
  \vspace{-1em}\begin{equation}
    \hat{x}=\frac{x'}{x}
  \vspace{-2em}\end{equation}
  \begin{itemize}
    \item Unobserved fundamentals $\{\bar u_i,\bar A_i,t_{kl}\}$ cancel out
    %\item They discipline the counterfactual through observed activity and traffic
  \end{itemize}
  \item A physical road $\{k,l\}$ has two directed lane allocations
  \vspace{-1em}\begin{equation}
    N_{kl}=n_{kl}+n_{lk}
  \vspace{-1em}\end{equation}
  and the policy shifts $\Delta_{kl}$ lanes from $l\to k$ to $k\to l$
  \vspace{-1em}\begin{equation}
    n'_{kl}=n_{kl}+\Delta_{kl},
    \qquad
    n'_{lk}=n_{lk}-\Delta_{kl}
  \vspace{-1em}\end{equation}
  \item The directed lane hats and mechanical cost hats are
  \vspace{-0.5em}\begin{equation}
    \hat n_{kl}=1+\frac{\Delta_{kl}}{n_{kl}},
    \qquad
    \hat n_{lk}=1-\frac{\Delta_{kl}}{n_{lk}},
    \qquad
    \hat c_{kl}^{\mathrm{mech}}=\hat n_{kl}^{\delta_1}
  \end{equation}
\end{itemize}
\end{frame}






\begin{frame}{Equilibrium Changes with Tidal Lanes}
\begin{itemize}
  \item For any vector of directed link-cost changes $\{\hat c_{kl}\}$, update the network by
  \vspace{-1em}\begin{equation}
    A'_{kl}=A_{kl}\hat c_{kl}^{-\theta}, \qquad B'=(I-A')^{-1}
  \vspace{-0.5em}\end{equation}
  \item (Link) traversal costs and (route) commuting costs change by
  \vspace{-0.5em}\begin{equation}\label{eq:chat}
    \hat c_{kl}=\hat\Xi_{kl}^{-\delta_1}\hat n_{kl}^{\delta_1}, \qquad \hat\tau_{ij} = \left( B'_{ij} / B_{ij} \right)^{-1/\theta}
  \vspace{-0.5em}\end{equation}
  \item Residence and workplace shares follow
  \vspace{-0.5em}\begin{equation}
    (\hat L_i^R)^{1+\beta} = \hat\lambda\sum_j s^R_{ij}\hat\tau_{ij}^{-\theta}(\hat L_j^F)^\alpha , \qquad 
(\hat L_i^F)^{1+\alpha} = \hat\lambda\sum_j s^F_{ji}\hat\tau_{ji}^{-\theta}(\hat L_j^R)^\beta
  \vspace{-0.5em}\end{equation}
  \item Directed traffic become
  \vspace{-0.5em}\begin{equation}
    \Xi'_{kl} = \sum_i\sum_j \kappa_{ij}^{kl\prime}L'_{ij}
 \end{equation}
\end{itemize}
\end{frame}

\begin{frame}{Counterfactual Algorithm}
\begin{enumerate}
  \item Choose treated physical roads and lane shifts $\{\Delta_{kl}\}$
  \item Compute directed lane hats $\{\hat n_{kl},\hat n_{lk}\}$
  \item Initialize traffic hats, for example $\hat\Xi_{kl}=1$
  \item Update link-cost hats by \eqref{eq:chat}
  \item Update bilateral costs using the matrix inverse and solve for $\hat L_i^R,\hat L_i^F$
  \item Recompute route intensities, directed traffic, and iterate to convergence
  \item Report mechanical, route-equilibrium, and full-equilibrium effects
\end{enumerate}
\end{frame}

\end{document}



%% 21
%\begin{frame}{Local Versus Global Congestion}
%\begin{itemize}
%  \item Local mechanical effect: hold routes and residence--workplace flows fixed.
%  \item Treated direction: lower traffic per lane and faster speed.
%  \item Donor direction: higher traffic per lane and slower speed.
%  \item Global equilibrium effect: commuters reoptimize routes and locations.
%  \item Network congestion: change in $2\sum_{k,l}d_{kl}m_{kl}\Xi_{kl}$ after reoptimization.
%\end{itemize}
%\end{frame}
%
%% 22
%\begin{frame}{Residence--Workplace Sorting}
%Tidal lanes can make directional commuting easier, especially for suburb-to-center work trips.
%\begin{equation}
%S
%=
%\frac{1}{2}\sum_i \left|l_i^R-l_i^F\right|.
%\end{equation}
%\begin{itemize}
%  \item If downtown becomes easier to reach from residential suburbs, workplace concentration may rise.
%  \item If suburbs become easier to access under the flipped return-trip asymmetry, residential decentralization may rise.
%  \item The model separates direct travel-time gains from induced changes in where people live and work.
%\end{itemize}
%\end{frame}
%
%% 24
%\begin{frame}{Data Needs}
%\begin{tabularx}{\textwidth}{@{}lX@{}}
%\toprule
%Object & Role in the tidal analysis\\
%\midrule
%Road geometry & Defines undirected roads $\{k,l\}$ and feasible directed links.\\
%Directional lanes & Determines $n_{kl}$ and the feasible tidal-lane reallocations.\\
%Directional traffic & Disciplines $\Xi_{kl}$ and $\Xi_{lk}$.\\
%Speeds or travel times & Converts traffic per lane into minutes and value-of-time costs.\\
%LODES or OD flows & Disciplines $L_{ij}$, residential shares, and workplace shares.\\
%Value of time & Converts static one-way travel-time changes into daily dollar congestion costs using a factor of two.\\
%\bottomrule
%\end{tabularx}
%\end{frame}
%
%% 25
%\begin{frame}{Calibration Choices}
%\begin{center}
%\begin{tabular}{@{}llll@{}}
%\toprule
%Parameter & Value & Interpretation & Source\\
%\midrule
%$\theta$ & 6.83 & commuting-cost sensitivity & AA / Ahlfeldt et al.\\
%$\alpha$ & 0.12 & productivity externality & AA / Ahlfeldt et al.\\
%$\beta$ & 0.10 & amenity externality & AA / Ahlfeldt et al.\\
%$\delta_1$ & $-1.1$ & speed-density elasticity & Barwick et al.; Yang et al.\\
%\bottomrule
%\end{tabular}
%\end{center}
%\begin{itemize}
%  \item Agglomeration parameters follow Allen--Arkolakis.
%  \item The congestion calibration uses Beijing evidence because it relates speed to traffic density per kilometer-lane.
%\end{itemize}
%\end{frame}
%
%% 26
%\begin{frame}{Why the Beijing Elasticity Is Useful}
%\begin{itemize}
%  \item Barwick et al. and Yang et al. estimate how speed changes with local traffic density.
%  \item Traffic density is converted to vehicles per kilometer-lane.
%  \item That object is close to the model's $\Xi_{kl}/n_{kl}$: directional link traffic per lane.
%  \item A speed elasticity of $-1.1$ implies that a 10\% increase in traffic per lane lowers speed by about 11\%.
%  \item Inverse speed and minutes per kilometer rise with elasticity $1.1$.
%\end{itemize}
%\end{frame}
%
%% 27
%\begin{frame}{Recovering Directional Flows}
%When direct directional traffic is unavailable, infer it using the model:
%\begin{equation}
%\Xi_{kl}=\sum_i\sum_j \kappa^{kl}_{ij} L_{ij},
%\qquad
%\Xi_{lk}=\sum_i\sum_j \kappa^{lk}_{ij} L_{ij}.
%\end{equation}
%\begin{itemize}
%  \item OD flows identify where commuters live and work.
%  \item Route intensities allocate OD commuters across network links.
%  \item Directional imbalance can be validated against directional counts or speed asymmetries.
%\end{itemize}
%\end{frame}
%
%% 29
%\begin{frame}{Decomposition Design}
%\begin{tabularx}{\textwidth}{@{}lX@{}}
%\toprule
%Scenario & Interpretation\\
%\midrule
%Mechanical & Shift lanes, hold OD flows and route shares fixed.\\
%Route equilibrium & Allow commuters to change routes, hold residence and workplace shares fixed.\\
%Full equilibrium & Allow routes, residences, and workplaces to adjust.\\
%\bottomrule
%\end{tabularx}
%\begin{itemize}
%  \item The first step isolates local engineering effects.
%  \item The second step measures network diversion and induced route demand.
%  \item The third step measures sorting and welfare effects.
%\end{itemize}
%\end{frame}
%
%% 30
%\begin{frame}{Expected Sorting Mechanism}
%\begin{itemize}
%  \item Tidal lanes reduce commuting costs in the dominant residence-to-work direction.
%  \item Lower suburb-to-center daily costs can make suburban residence and central employment more attractive jointly.
%  \item The same policy can reduce peak bottleneck delays while increasing residence--workplace specialization.
%  \item The welfare effect depends on whether time savings dominate induced congestion and any sorting externalities.
%\end{itemize}
%\end{frame}
%
%% 31
%\begin{frame}{Planned Empirical Result Slides}
%\begin{itemize}
%  \item Map of baseline directional traffic imbalance.
%  \item Map of roads eligible for a one-lane tidal reallocation.
%  \item Local congestion effects for treated directions.
%  \item Global network congestion effects after route reoptimization.
%  \item Residence and workplace reallocation maps.
%  \item Welfare and value-of-time congestion-cost tables.
%\end{itemize}
%\end{frame}
%
%% 32
%\begin{frame}{Limitations and Assumptions}
%\begin{tabularx}{\textwidth}{@{}lXX@{}}
%\toprule
%Assumption & Why it is useful & Caveat\\
%\midrule
%Directed traffic as commuting traffic & Aligns congestion with work trips & Non-work travel and mode choice may matter.\\
%Symmetric physical road & Separates infrastructure from directional demand & Some roads may have fixed asymmetric design.\\
%Log-linear speed-density curve & Uses transparent literature calibration & May miss capacity thresholds and queuing.\\
%One-lane reallocation & Matches tidal-lane policy design & Requires feasible reversible-lane operations.\\
%Static daily equilibrium & Keeps sorting and route choice tractable & Does not model departure-time choice.\\
%\bottomrule
%\end{tabularx}
%\end{frame}
%
%% 33
%\begin{frame}{Takeaways}
%\begin{enumerate}
%  \item Tidal commuting requires directed traffic even when physical roads are symmetric.
%  \item The residence--workplace matrix determines the congested direction in the static network.
%  \item Tidal lanes hold total road capacity fixed and reallocate lanes toward the more congested direction.
%  \item Evaluation should separate local mechanical congestion effects, global route-equilibrium effects, and full residence--workplace sorting effects.
%  \item Value-of-time congestion costs are twice the static one-way travel-time cost because the return-trip asymmetry flips.
%\end{enumerate}
%\end{frame}
